% Generated draft: requires mathematical and source review.
\documentclass[11pt]{article}
\usepackage[margin=25mm]{geometry}
\usepackage{fontspec}
\setmainfont{Latin Modern Roman}
\usepackage{amsmath,amssymb,mathtools,unicode-math}
\setmathfont{Latin Modern Math}
\usepackage{enumitem,needspace,xurl,hyperref}
\setlength{\parindent}{0pt}
\setlength{\parskip}{5pt}
\setlength{\emergencystretch}{3em}
\newcommand{\N}{\mathbb{N}}
\newcommand{\Z}{\mathbb{Z}}
\newcommand{\Q}{\mathbb{Q}}
\newcommand{\R}{\mathbb{R}}
\newcommand{\C}{\mathbb{C}}
\newcommand{\F}{\mathbb{F}}
\newcommand{\A}{\mathbb{A}}
\newcommand{\PP}{\mathbb{P}}
\newcommand{\E}{\mathbb{E}}
\newcommand{\eps}{\varepsilon}
\newcommand{\dd}{\,\mathrm{d}}
\begin{document}
\section*{UnsolvedMath ID: 20003109}
\textbf{Source problem code:} catalogue entry
\subsection{Definitions and mathematical statement}
Problem 3.6, attributed to F. Gu\'eritaud, asks for the cone of tangent directions at $X\in\mathcal T(S)$ shortening the lengths of all simple closed curves. The supplied source verification places this in reduced Teichm\"uller space for a finite-type surface with nonempty boundary. Boundary lengths and the class of curves being tested must be specified; fixing a boundary length would exclude strict shortening of that boundary curve if it were included. The following arc-metric/doubling paragraph in the extraction belongs to the context of Problem 3.7, not to this problem.

Let $\mathscr C$ be the intended set of simple closed curves and write $\ell_\alpha=\ell_\alpha(X)$. Define the literal shortening cone
\[
 L_X=\{v:d\ell_\alpha|_X(v)<0\text{ for every }\alpha\in\mathscr C\}.
\]
In the compact-dual framework of the aid, let $\mathscr K_X$ be the compact closure of the normalized curve laminations $\alpha/\ell_\alpha$ in the relevant measured-lamination/current model. Length derivative extends continuously and linearly to it; put $f_v(\lambda)=d\ell_\lambda|_X(v)$. Define
\[
 U_X=\{v:\max_{\lambda\in\mathscr K_X}f_v(\lambda)<0\},
 \qquad
 W_X=\{v:\max_{\lambda\in\mathscr K_X}f_v(\lambda)\leq0\}.
\]
The first is uniform shortening, equivalently $d\ell_\alpha(v)/\ell_\alpha\leq-c$ for some $c>0$ and every selected curve. The aid reports
\[
 U_X\subset L_X\subset W_X,\qquad
 \overline{U_X}=\overline{L_X}=W_X,
\]
with the closure argument requiring the nonempty uniform cone in the intended setting. Its residual decomposition is
\[
 L_X=U_X\ \sqcup\ R_X,\qquad
 R_X=\{v\in L_X:\max_{\mathscr K_X}f_v=0\}.
\]
A zero supporting lamination for $v\in R_X$ cannot be a nonzero finite weighted multicurve, since every selected simple-curve derivative is strictly negative.

For a vector already satisfying all literal strict inequalities, the aid sharpens the residual test to
\[
 v\in R_X\quad\Longleftrightarrow\quad
 \exists\text{ pairwise distinct }\alpha_j\in\mathscr C:
 \ \ell_{\alpha_j}(X)\to\infty,\quad
 \frac{d\ell_{\alpha_j}|_X(v)}{\ell_{\alpha_j}(X)}\to0\text{ from below}.
\]
Every projective accumulation of this sequence is a non-rational zero supporting lamination. The sequence condition alone does not replace the requirement that all simple curves are shortened.

The aid identifies the uniform part with the antipode of the DGK strip-deformation cone parameterized by the pruned filling arc complex. That identification requires the DGK surface, boundary and tested-geodesic hypotheses to match the chosen $\mathscr C$; they are not fully stated here. For the one-holed torus it reports Gu\'eritaud's description of the literal cone as the antipode of the closed weak polygon with all rational sides removed, understood in the theorem's projective or affine normalization and excluding the cone apex. No precise polygon or general-boundary classification is supplied by the summary. These descriptions must not collapse literal and uniform shortening into the same condition.
\subsection{Short English statement}
Describe directions that strictly shorten every simple closed curve on a bordered surface. The supplied aid separates uniform shortening from a possible boundary residual. A residual vector still shortens every curve, but a sequence of longer and longer curves has normalized derivatives approaching zero. The source's adjacent arc-metric paragraph is not part of this question.
\subsection{Sources}
\begin{itemize}
\item Local UnsolvedMath registry, permanent ID $20003109$, source code catalogue entry; canonical location \texttt{aim-topology-notes.json}, \texttt{notes[196]}.
\item AIM \emph{Problems on the Lipschitz metric on Teichm\"uller space}, Problem 3.6 (F. Gu\'eritaud), \url{https://aimath.org/pastworkshops/lipschitzteichproblems.pdf}. The supplied verification locates reduced bordered Teichm\"uller space on page 3, lines 123--125, this problem on lines 126--127, separator on line 128, and the next problem's arc/doubling context on lines 129--135. It marks the corrected statement verified; no PDF was accessed here.
\item DGK strip-cone parameterization and Gu\'eritaud's one-holed-torus theorem, as named in the research summary, without full bibliographic references.
\item Research attempt 2 is marked partial result, proved, and low novelty confidence, with an explicit machine-generated, non-peer-reviewed research-aid warning requiring independent verification before citation.
\end{itemize}
\end{document}
